\chapter{Conclusions sur la construction générative et ses réalisations}
\chaptermark{Conclusions sur la construction générative}
\label{hmt:sistematica:conclusiones}

Le développement établit une relation constructive entre arithmétique modulaire, dynamique orientée, génération de valeurs, incidence exceptionnelle et récupération de l’information. L’unité de cette relation réside dans l’état enrichi produit par APP--TRIT--TPK. Chaque réalisation reçoit une partie déterminée de sa structure et conserve les données exigées par l’opération suivante. L’évaluation numérique, la forme géométrique et le registre d’une exécution sont ainsi articulés par des applications explicites.

Les conclusions principales concernent trois questions mathématiques indissociables : comment une valeur est produite, quelle structure sa production conserve et sous quelles opérations elle peut être récupérée. Les réalisations de la Mécanique Dimensionnelle ajoutent une quatrième détermination : comment cette structure s’exprime en observables avec des unités et des conditions de couplage déclarées.

\section{Restitution arithmétique et mémoire des opérations}

La construction initiale d’APP fixe deux niveaux : la position sur la grille et l’évaluation arithmétique qui y est effectuée. La réduction nonadique conserve exactement l’entier lorsqu’elle retient le reste et le quotient. La réduction ternaire équilibrée possède la même propriété avec son représentant orienté et son quotient de prolongement. Les identités de composition démontrent que les deux registres peuvent être transportés entre les deux décompositions.

En particulier, la restitution \(n=\rho_9(n)+9q_9(n)\) et la compatibilité \(c_3(n)=c_3(\rho_9(n))+3q_9(n)\) permettent de suivre l’information entière lors des changements de représentation. Les retenues de l’addition et les termes croisés du produit sont calculés sur ces coordonnées. L’information récupérable possède, dès cette étape, un contenu algébrique déterminé \ref{thm:app-reconstruccion}, \ref{trit:prop:operaciones}.

La composition des trajectoires ajoute une mémoire ordonnée. Si une transition agit par \(U_\varepsilon(m)=C_\varepsilon m+
\kappa_\varepsilon\), la concaténation satisfait
\[
 C_{\beta\alpha}=C_\beta C_\alpha,
 \qquad
 \kappa_{\beta\alpha}=C_\beta\kappa_\alpha+
 \kappa_\beta.
\]
La seconde trajectoire transporte la mémoire accumulée avant d’ajouter son incrément. Cette loi de cocycle explique pourquoi le registre chronologique contient davantage d’information qu’une retenue agrégée. Sa démonstration fixe l’opération à conserver lors de la composition des réalisations \ref{x_reciprocidad:iv:cont:concatenacion}.

Le calendrier nonadique apporte une autre distinction exacte. Le retour de neuf phases incrémente la coordonnée de tour. L’extension \(C_{12}\to C_{108}\to C_9\), avec sa retenue, exprime la compatibilité entre la phase visible et le registre dodécaphasique. La continuation enrichie conserve en outre les tours successifs et les incidences. La périodicité d’une lecture et l’évolution de sa mémoire constituent donc deux propriétés simultanées du même transport.

Le transfert entre les complétions quadratique et cubique fournit une autre identité exacte de conservation. Dans le cas nonadique,
\[
 (72,27)\longmapsto(73,27)
 \longmapsto(10\cdot73-1,10\cdot27+1)=(729,271).
\]
La première opération incorpore l’unité enregistrée ; la seconde change l’échelle et redistribue une unité par une correction de somme nulle. La somme normalisée reste égale à un. Cette construction donne un contenu précis à l’extrême et moyenne raison holographique et à son prolongement conservatif entre échelles \ref{x_reciprocidad:lib03:prop-completacion}.

\section{Détermination des régions et de la quadrirelation coinductive}

L’énumération finie identifie les configurations, émissions et mots admissibles sur lesquels agissent les relèvements. La chaîne \(w_6\to w_{12}\to w_{18}\to w_{24}\to w_{30}\to R_{36}\to G_9\) conserve les préfixes et les conditions de transport lors du prolongement des réalisations régionales. La compatibilité cylindrique et la contraction assurent le passage à une profondeur arbitraire.

Les sorties régionales \(\pi_{\rm HMT}\), \(\varphi_{\rm HMT}\) et \(e_{\rm HMT}\) réunissent un développement positionnel et une provenance opératoire. La classification des cinq régions de \(\pi\), l’incidence de l’axe de \(\varphi\), les continuations conjointes et l’involution spéculaire conservent leurs distinctions à la limite. Les identifications analytiques ultérieures certifient des propriétés de ces sorties et permettent de les comparer dans des langages mathématiques établis.

La connexion nonadique relie la génération régionale à la structure du retour. L’incidence modale provient du sélecteur tritique et de sa suite d’opérations. Ses puissances engendrent des récurrences ; les marques régionales sont ensuite appliquées au support des retours. Cette antériorité distingue l’émetteur régional, la récurrence d’une réalisation modale et l’évaluation d’un rapport de frontière. La composition permet de calculer leurs relations sans identifier des objets de types différents.

Le registre \(K\) est déterminé à partir des bandes régionales. La première avancée nulle de l’horloge de conversion localise l’antécédent nonadique ; les défauts orientés et le lecteur phase--charge produisent le descripteur. Le répertoire terminal, les cylindres ternaires et les incidences exceptionnelles délimitent les candidats. La condition d’hétérotypie sélectionne le registre dans ce domaine fini. L’inversion ultérieure récupère son ordre, ses fenêtres et sa réalisation signée. L’unicité conserve explicitement le domaine des candidats et les conditions du sélecteur.

La clôture entière de structure fine compose les lectures régionales avec la coordonnée rationnelle de \(K\). Sa retenue compatible détermine une unique section infinie, et l’identité des queues distingue la frontière d’une fenêtre isolée de la restriction de la limite. Dans les notations de son développement,
\[
 \alpha_{\rm HMT}
 =\pi_{\rm HMT}+e_{\rm HMT}-\varphi_{\rm HMT}
  -4-\kappa_{\rm per}.
\]
L’égalité conserve l’origine de tous ses termes et la normalisation décimale complète \eqref{x_reciprocidad:eq:alpha-salida-a-completa}.

La représentation analytique développe cette même précoordonnée au moyen d’une fonction définie à tous les degrés. Si \(a\) est la précoordonnée engendrée et \(q=1/729\), le coefficient de liaison
\[
 \lambda=-\frac{E_9(a)(1-qa^3)}{a^{10}}
\]
détermine le prolongement géométrique \(E_\infty(X)=E_9(X)+\lambda X^{10}/(1-X^3/729)\). Les coefficients ultérieurs satisfont \(b_{10+3n}=q^n\lambda\) et \(b_{11+3n}=b_{12+3n}=0\). L’unicité est démontrée dans cette classe de prolongements, et la racine est identifiée à la coordonnée commune. La seconde voie apporte une représentation analytique corrélée avec contrôle de convergence et compatibilité projective \ref{x_reciprocidad:subsec:alpha-lector-completo-rev08}.

\section{Réalisation archimédienne et incidence exceptionnelle du continu}

Les systèmes de troncature conservent les trajectoires admissibles et leurs registres. La prolongeabilité démontrée garantit l’existence de sections compatibles. Les partitions cylindriques fournissent une évaluation archimédienne ; la structure d’incidence demeure dans la projection conjuguée. La valeur réelle est déterminée par la contraction, tandis que l’orientation, la frontière et la mémoire conservent les distinctions de l’état.

Cette double projection organise conjointement le transport spectral, le bilan solénoïdal, l’incidence de jauge, la cohérence cohomologique et la droite déterminante. Les identités de naturalité permettent d’effectuer les raffinements et les réalisations dans des ordres compatibles. Le continu discret reçoit ainsi une structure opératoire commune sur laquelle agissent ses évaluations ultérieures.

L’incidence exceptionnelle se développe selon une suite de constructions concrètes. Les codes et les marquages déterminent des recollements réticulaires ; le voisin de Leech est démontré positif, pair, unimodulaire et sans racines dans la métrique et la chambre déclarées. L’identification réticulaire utilise conjointement ces propriétés. La construction des oscillateurs et de l’algèbre réticulaire tordue, suivie du secteur de parité de l’orbifold, détermine la réalisation de \(V^\natural\) et son action du Monstre~\ref{x_reciprocidad:exc:moonshine}.

L’importance de cette suite tient à la généalogie liée de la valeur et de la symétrie et à leurs objets d’action précis. La symétrie exceptionnelle agit sur le code, le réseau ou l’algèbre correspondants. Les réductions dimensionnelles et les réalisations de dualité reçoivent ces supports avec leurs conditions. Le raffinement numérique et le prolongement exceptionnel conservent ainsi leurs applications de liaison dans l’état commun.

\section{Conservation bilatérale et réversibilité à profondeur arbitraire}

La loi bilatérale fournit une conclusion opératoire indépendante de la dimension. Pour deux isométries \(\mathcal B,\mathcal C:H\to H'\) et \(a>0\), on définit \(Q_-=a\mathcal B-\mathcal C\) et \(Q_+=(a+1)\mathcal B+\mathcal C\). L’identité
\[
 (a+1)Q_-^*Q_-+aQ_+^*Q_+
 =\bigl((a+1)^3+a^3\bigr)I
\]
résulte de l’annulation des termes croisés. Sa validité comprend les espaces de Hilbert de dimension infinie et les transports entre supports distincts ayant un domaine et un codomaine communs.

Dans la spécialisation nonadique, \(a=8\), la normalisation se factorise au moyen des poids \(72/73\) et \(1/73\), suivis d’une rotation des deux canaux. L’isométrie obtenue possède un inverse à gauche de norme un. La récupération conserve ainsi un contrôle optimal de l’erreur conjointe dans la norme déclarée. Le facteur \(73=72+1\) intervient dans cette identité ; ses autres réalisations comme norme ou indice réticulaire conservent leurs propres démonstrations.

Le transport des observables précise le contenu de la conservation. Pour un lecteur \(G\) et un transport relatif unitaire \(R\), la lecture induite est
\[
 \mathcal T_a(G)=
 \frac{a(a+1)G+R^*GR}{a(a+1)+1}.
\]
L’égalité \(\mathcal T_a(G)=G\) équivaut à la commutation \(GR=RG\). On obtient ainsi un critère de symétrie relative distinguant les observables invariantes de celles qui changent pendant la transduction \ref{x_reciprocidad:lib04:teo-lector}.

La règle d’enregistrement à profondeur arbitraire conserve à chaque étape la composante prolongée et son complément. Le bilan télescopique produit une isométrie finie et, sous la borne de contraction démontrée, une isométrie à registre infini. Pour \(a=8\), le taux uniforme est \(r_8=729/1241<1\). Les \(N\) premiers compléments permettent de récupérer l’état avec une erreur de troncature bornée par \(r_8^N\|v\|\) ; une erreur conjointe \(\varepsilon\) ajoute exactement \(\varepsilon\) à cette borne. Le registre complet admet l’inverse convergent du théorème~\ref{x_reciprocidad:lib04:teo-archivo}.

La conservation évolutive de l’unité s’exprime donc par des identités locales, des compositions isométriques et une récupération à la limite. La modification de la forme ou du support est compatible avec la restitution de l’état lorsque les canaux et les incidences spécifiés par le théorème sont conservés.

\section{Réciprocité géométrique et conservation extérieure}

Le registre de couplage produit une direction orthogonale au mode uniforme. Son flot diagonal positif a pour déterminant un. Les rayons varient de manière compensée ; les mineurs de Plücker conservent leur signe et leur nullité, et les quotients équilibrés conservent leur valeur. La même direction engendre sur les plans réticulaires une déformation de covolume unité.

La réciprocité de cette déformation s’exprime par \(\Lambda_t^*=\Lambda_{-t}\). La symétrie ternaire commutante persiste le long du flot, tandis que l’intégralité et la réalisation exceptionnelle conservent le domaine spécifique où elles ont été démontrées. Le théorème~\ref{x_reciprocidad:rec:thm:red} détermine simultanément ces propriétés.

Dans la carte \(s=(1-z)^{-1}\), l’inversion radiale \(z\mapsto1/\overline z\) devient \(s\mapsto1-\overline s\). Le cercle unité correspond à la droite \(\Re s=1/2\), le point singulier étant exclu. La réciprocité thêta--Mellin, la structure polaire et les réalisations de Poincaré et de Klein conservent les conditions analytiques et géométriques de leurs propres lecteurs. L’orientation enregistrée permet d’inverser la lecture radiale là où la lecture scalaire réunit des réalisations conjuguées.

Ces identités élargissent la signification de la constante de couplage : le registre arithmétique détermine un repère récupérable, une direction de déformation et des transports conservant des invariants extérieurs. La correspondance entre les réalisations s’obtient par composition des applications construites.

\section{Action, géométrie angulaire et observables matérielles}

La réalisation de Mécanique Dimensionnelle conserve une antériorité explicite. L’échelle d’action reçoit la coordonnée de structure fine et la lecture régionale de \(\varphi\) ; la norme locale d’ordre seize fixe sa valuation décimale. Les sections aller et retour de Planck conservent leur quotient \(R_{\rm act}\). Le couple angulaire et ses canaux conjugués sont ensuite évalués avec leur orientation et leurs unités~\ref{ii:sec:ii-definicion-planck}.

L’inclusion marquée d’une hexade dans une octade produit les incidences de cardinaux \(90\) et \(120\). La chaîne dodécaphasique détermine les multiplicités de la fonctionnelle de Barbero--Immirzi et le coefficient de pôle \(1/24\). Son évaluation utilise l’angle, le contre-angle et les canaux précédemment construits. La parité sous inversion angulaire conserve l’orientation de cette lecture~\ref{ii:sec:barbero}.

Le transport angulaire vers l’incidence de Witt fournit une composition explicite avec le lecteur de mélange. Son inverse à gauche et son identité de Gram conservent l’information angulaire transportée \ref{ii:thm:ii-witt-angular-transport}. La structure constitutive électrique et magnétique reçoit également ses modes, ses rapports de retour et ses échelles ; ses identités permettent de reconstruire les coordonnées angulaires dans le domaine spécifié.

La section gravitationnelle contragrédiente de l’action satisfait
\[
 G_a=\frac{c_{\rm int}^{3}R_{108}^{2}}{\hbar_a},
 \qquad
 \frac{G_a\hbar_a}{c_{\rm int}^{3}}=R_{108}^{2}.
\]
La condition de coïncidence des rayons fixe cette section dans la classe déclarée. L’inversion du rapport d’action compense le rapport gravitationnel et conserve l’unité aréale \ref{ii:sec:ii-definicion-gravedad}. La définition de Boltzmann relie ensuite l’énergie par décision à la section thermique. Pour l’horloge de cent huit événements, \(k_B\Theta_{{\rm clk},a}=h_a/(108t_0\log3)\). L’individualisation du transducteur conserve la section thermique et les bases d’unités utilisées \ref{ii:sec:ii-definicion-boltzmann}.

Les observables matérielles proviennent de la persistance des extensions et de leurs registres de parcours. Le centre électronique possède une sélection ponctuelle déterminée par l’atlas ; les autres classes utilisent les opérateurs de leurs fibres, orbites ou multisections. La signature entière, le caractère et les tours transmettent cette information à la réalisation de masse. La comparaison métrologique porte sur les évaluations et les couplages spécifiés, en conservant la distinction entre opérateur, spectre et scalaire individualisé.

\section{Reconstruction et classification cardinale du continu engendré}

La construction des cylindres démontre une couverture des valeurs avec récupération des registres. La clôture cardinale ajoute une classification de ces valeurs dans la clôture généalogique \(\mathfrak G\) déterminée par sa sémantique des termes et sa définissabilité. La dénombrabilité interne des niveaux et le lemme de condensation permettent d’assigner à chaque coupure rationnelle sa première étape \(\beta(r)\) et son premier indice \(n(r)\) dans l’énumération de cette étape.

L'application
\[
 j_{\rm HMT}(r)=\omega\beta(r)+n(r)
\]
est injective sur le domaine de la droite engendrée \(\mathbb R_{\rm HMT}^{\mathfrak G}\). Le codage des bons ordres dénombrables construit l’injection de \(\omega_1^{\mathfrak G}\) dans cette même droite. Cantor--Bernstein conclut
\[
 \left|\mathbb R_{\rm HMT}^{\mathfrak G}\right|
 =\aleph_1^{\mathfrak G},
 \qquad
 \mathfrak G_{\rm HMT}(h,m_\infty)
 \models\mathrm{CH}.
\]
Il s’agit de la clôture affirmative de l’hypothèse du continu établie dans le théorème~\ref{ix:thm:decision-continuo-hmt}.

La conclusion comprend tous les sous-ensembles internes de cette droite : chacun est dénombrable ou possède le cardinal du continu engendré. La représentation \(\mathfrak G_{\rm HMT}(h,m_\infty)=L[h,m_\infty]\) identifie le domaine de la construction après le développement de sa généalogie. Les transferts vers un domaine ensembliste différent conservent comme hypothèses la couverture et la coïncidence des objets sur lesquels agissent les injections. La formulation maintient ainsi ensemble le résultat cardinal, sa sémantique et ses quantificateurs.

Les réalisations terminales spectrale, solénoïdale, de jauge, cohomologique et déterminantale s’articulent à travers la même structure discrète. Leurs théorèmes spécifient l’opérateur, la positivité, la coercivité ou la descente utilisés par chaque réalisation, avec les identités de naturalité correspondantes. La factorisation conjointe du chapitre~\ref{chap:terminal-arquitectura-natural-comun} conserve cette origine et le domaine de chaque conclusion.

\Needspace{27\baselineskip}
\section{Contenu démonstratif et portée de l’intégration}

La force démonstrative du travail provient de la composition de résultats de types différents. Les identités symboliques démontrent des égalités sur leurs domaines quantifiés. Les énumérations exhaustives résolvent des sélecteurs et des dénombrements finis. Les bornes rationnelles contrôlent les approximations et les racines. Les arguments de compatibilité, de convergence et de naturalité justifient les prolongements et les passages à la limite.

Les formalisations en Lean vérifient les déclarations précises qu’elles contiennent, avec leurs types, prémisses et dépendances. Les calculs exécutables certifient les instances ou domaines finis définis par leurs algorithmes. La correspondance entre énoncé mathématique et certificat permet d’assigner à chaque vérification sa portée précise. La réalisation physique conserve en outre la base dimensionnelle et la règle de couplage avec l’observable.

Le résultat d’ensemble est une généalogie opératoire : l’arithmétique discrète produit des états orientés ; leurs prolongements déterminent des constantes et un continu doté d’incidences ; leurs registres admettent des transformations conservatives et une récupération ; leurs réalisations dimensionnelles fournissent des observables structurées. L’unité du développement s’exprime dans ces relations démontrées. La valeur, la symétrie et l’information récupérable sont liées par la construction qui les produit.
